Change-point tests based on persistent homology use sliding windows to extract local structural features, but dependence induced by shared observations can distort random-permutation reference distributions. This study examines the finite-sample calibration of a test combining a fixed vector representation of persistence diagrams with a maximum scan statistic. A Gaussian window-mean counterexample first shows that independent and identically distributed raw observations do not ensure exchangeability of the derived window sequence. Naive unrestricted window-level permutation and a dependent multiplier bootstrap are then evaluated under fixed-window-count and fixed-observation-length designs across 28 null configurations, with mean--covariance features as a non-topological comparator. At the nominal level of 0.05, all 20 overlapping-window configurations are classified as materially anti-conservative for each representation; the empirical rejection rate for the persistent-homology representation reaches 0.9756. The multiplier procedure lowers rejection rates in every matched configuration, but all eight non-overlapping configurations are materially conservative for each representation. The effect of serial dependence on the stride contrast varies across generators and representations. Under a circle-to-figure-eight alternative with matched first and second marginal moments, size-adjusted powers are 1.0000 and 0.0620 for the persistent-homology and moment representations, respectively. A mean-shift negative control reverses this ordering, consistently with the translation invariance of distance-based persistent homology. These results demonstrate condition-dependent calibration for the studied implementation: lower rejection rates do not necessarily imply nominal calibration, and representation sensitivity and Type-I error control require separate evaluation.

\noindent\textbf{Keywords}: persistent homology; change-point detection; overlapping sliding windows; exchangeability; finite-sample calibration.
